\scanpage{177}
\mspage{93}
% archive sheet number: 176
% lines 1-6
\uline{Lemme} 4.4.12. Soit \deleted{\unsure{les}} $(C_j)_{j\in J}$ une famille de \inserted{petites} catégories grandes devant $\pi$, indexée par un \inserted{petit} ensemble $J$ \deleted{tel que $\operatorname{card}J\leq\pi$}. Alors $\prod C_j$ est grande devant $\pi$, et les projections (pour $j_0\in J$) $\prod C_j\longrightarrow C_{j_0}$ sont cofinales.

% lines 7-16
Prouvons d'abord la cofinalité, qui est la plus importante pour nous. C'est \unsure{évident} si $J=\{j_0\}$, sinon on a $J=J'\cup\{j_0\}$, et $\prod_{j\in J}C_j\simeq C'\times C_{j_0}$, où $C'=\prod_{j\in J'}C_j$, donc il faut voir que $C'\times C_{j_0}\longrightarrow C_{j_0}$ est cofinal, i.e. que les catégories ${}_{\xi}\backslash(C'\times C_{j_0})$ (pour $\xi\in\operatorname{Ob}C_{j_0}$) sont $0$-connexes, or \deleted{elles sont} elles se \unsure{ramènent|réduisent}: $C'\times{}_{\xi}\backslash C_{j_0}$, et ${}_{\xi}\backslash C_{j_0}$ est $0$-connexe, il reste à voir que $C'$ est $0$-connexe.

% lines 17-24
Or \unsure{faire|fera} attention à cet égard qu'un produit \uline{infini} de catégories $0$-connexes n'est pas forcément $0$-connexe, donc \deleted{\unsure{il}} la \deleted{elle n'est pas} \inserted{cofinalité} n'est pas totalement tautologique! Mais si on prouve que la catégorie $\prod_{j\in C}C_j$ est
